\chapter{Comprendre avant de lancer}

Ce chapitre ne contient aucune commande. Il contient les cinq idées sans
lesquelles toutes les commandes des chapitres suivants ressembleraient à de la
magie.

\section{Qu'est-ce qu'un modèle, au juste ?}

Imaginez un très grand tableau de nombres. Pas quelques centaines : sept
milliards. Ces nombres s'appellent des \textbf{paramètres}.

Quand vous écrivez une phrase au modèle, elle traverse ce tableau. À chaque
étape, des multiplications et des additions transforment votre phrase en une
autre suite de nombres, jusqu'à produire une réponse : quel mot vient ensuite ?

C'est tout. Un modèle de langue est une immense machine à deviner le mot
suivant. Sa culture, son style, sa capacité à raisonner — tout est rangé dans la
valeur de ces sept milliards de nombres.

\begin{nkretenir}
Un modèle, c'est un tableau de nombres. « Entraîner » veut dire : \emph{modifier
ces nombres} pour que les réponses ressemblent davantage à celles qu'on
souhaite.
\end{nkretenir}

\section{Qu'est-ce qu'entraîner ?}

Prenons une image. Vous apprenez à un enfant à reconnaître un chat. Vous montrez
une photo, il dit « chien ». Vous corrigez : « non, chat ». Il ajuste sa manière
de voir. Vingt photos plus tard, il ne se trompe plus.

L'entraînement d'un modèle suit exactement ce schéma, avec trois différences :

\begin{enumerate}
  \item On ne corrige pas « un peu », on mesure \textbf{de combien} il s'est
        trompé. Ce nombre s'appelle la \textbf{perte} (\emph{loss} en anglais).
        Plus il est bas, mieux c'est.
  \item On ne corrige pas au hasard : on calcule, pour \emph{chacun} des
        paramètres, dans quel sens le pousser pour que la perte diminue. Cet
        ensemble de directions s'appelle le \textbf{gradient}.
  \item On recommence des milliers de fois.
\end{enumerate}

Un \textbf{pas} d'entraînement, c'est donc : montrer un exemple, mesurer la
perte, calculer le gradient, pousser un peu les paramètres dans le bon sens.

\begin{nkretenir}
La perte descend $=$ le modèle apprend. C'est le seul chiffre à surveiller
pendant les premières heures.
\end{nkretenir}

\section{Pourquoi on ne réentraîne pas tout : l'idée du LoRA}

Modifier sept milliards de nombres demande une carte graphique à plusieurs
dizaines de milliers d'euros. Ce n'est pas votre cas, ni le mien.

L'astuce s'appelle \textbf{LoRA} — \emph{Low-Rank Adaptation}, « adaptation de
rang faible ». Elle tient en une phrase :

\begin{quote}
On \textbf{gèle} les sept milliards de paramètres d'origine, et on ajoute à côté
un tout petit nombre de paramètres neufs, qui sont les seuls à apprendre.
\end{quote}

Reprenons l'image de l'enfant. Plutôt que de lui refaire toute son éducation, on
lui donne une paire de lunettes. Les lunettes sont minuscules comparées à
l'enfant, mais elles suffisent à changer ce qu'il voit.

Les chiffres, sur le modèle utilisé dans ce cours :

\begin{center}
\begin{tabular}{lr}
\toprule
\textbf{Ce qui existe} & \\
Paramètres du modèle d'origine (gelés) & $\approx 7\,600\,000\,000$ \\
\midrule
\textbf{Ce qu'on ajoute et qui apprend} & \\
Paramètres LoRA (rang 8) & $20\,185\,088$ \\
Proportion & $\mathbf{0{,}27~\%}$ \\
\bottomrule
\end{tabular}
\end{center}

On entraîne donc moins de trois millièmes du modèle. Et ça marche : c'est ce que
montrent les mesures du chapitre~4.

\begin{nkretenir}
LoRA : le modèle d'origine ne bouge pas, on lui greffe de petites pièces qui,
elles, apprennent. C'est ce qui rend l'entraînement possible sur une carte de
8~Go.
\end{nkretenir}

\section{Les deux voies : partir d'un modèle, ou partir de rien}

\subsection*{Voie A — partir d'un modèle existant}

Vous prenez un modèle déjà entraîné par quelqu'un d'autre — il sait déjà lire,
écrire, raisonner — et vous lui apprenez \emph{votre} domaine. C'est la voie
décrite aux chapitres 3 à 7.

\textbf{Coût :} quelques heures à quelques jours. \textbf{Résultat :} un modèle
qui parle déjà bien et qui connaît en plus votre sujet.

\subsection*{Voie B — partir de zéro}

Vous créez un modèle vide, avec des nombres tirés au hasard, et vous lui
apprenez tout : le vocabulaire, la grammaire, le sens. C'est la voie du
chapitre~8.

\textbf{Coût :} des mois et des machines que vous n'avez pas, pour un modèle de
grande taille. \textbf{Mais} c'est parfaitement faisable — et instructif — pour
un petit modèle sur une tâche étroite.

\begin{nkpiege}
Beaucoup de débutants veulent « créer leur propre IA de zéro ». C'est
compréhensible et c'est presque toujours une erreur de départ : vous
dépenseriez des mois à réapprendre au modèle ce que la langue française est
déjà. \textbf{Commencez par la voie A.} Le chapitre~8 vous montrera la voie B
sur un problème à sa mesure.
\end{nkpiege}

\section{Généraliste ou spécialisé ?}

Un modèle \textbf{généraliste} répond correctement à tout, sans exceller nulle
part. Un modèle \textbf{spécialisé} connaît un domaine en profondeur.

La différence ne tient pas à la méthode — elle est identique — mais au
\textbf{corpus} : le texte dont il apprend. Un corpus de conversations variées
donne un généraliste ; un corpus de documentation médicale donne un modèle
médical.

\begin{nkretenir}
Vous ne choisissez pas « généraliste » ou « spécialisé » par un réglage. Vous le
choisissez en choisissant vos textes. Le chapitre~3 est donc le plus important
du cours, même s'il ne contient presque pas de commandes.
\end{nkretenir}

\section{Le vocabulaire, en une page}

\begin{center}
\begin{longtable}{p{3.6cm}p{10.4cm}}
\toprule
\textbf{Mot} & \textbf{Ce que ça veut dire} \\
\midrule
\endhead
Paramètre & Un des nombres du modèle. \\
Perte (\emph{loss}) & De combien le modèle s'est trompé. Plus bas = mieux. \\
Gradient & Dans quel sens pousser chaque paramètre pour baisser la perte. \\
Pas (\emph{step}) & Un exemple montré $+$ une correction appliquée. \\
Époque (\emph{epoch}) & Un passage complet sur tout le corpus. \\
Corpus & L'ensemble des textes dont le modèle apprend. \\
Token & Un morceau de mot. « bonjour » peut valoir 2 tokens. \\
LoRA & Les petites pièces ajoutées qui apprennent à la place du modèle. \\
Rang & Taille de ces pièces. Plus grand = apprend plus, coûte plus. \\
Adam & La méthode qui décide de l'ampleur de chaque correction. \\
Checkpoint & Une photo de l'entraînement, pour pouvoir reprendre après une coupure. \\
Quantification & Ranger les nombres sur moins de bits pour tenir en mémoire. \\
VRAM & La mémoire de la carte graphique. C'est elle qui vous limite. \\
\bottomrule
\end{longtable}
\end{center}

\begin{nkexo}
Sans relire le tableau, expliquez à voix haute — vraiment à voix haute — la
différence entre un \emph{pas} et une \emph{époque}, puis pourquoi on gèle le
modèle d'origine dans LoRA. Si vous butez, relisez la section concernée : les
chapitres suivants supposent ces deux idées acquises.
\end{nkexo}
